% ===========================================================================
%  附录：补充分析与讨论
%  format2026.doc 第四条规定正文不超过 30 页、附录页数不限。正文篇幅已达上限，
%  故把稳健性、灵敏度、成因分解与适用边界这类"支撑性"讨论移入本附录。所有数据、
%  公式与结论一条未删；正文对应位置保留了结论摘要（含全部关键数字）并指向本节。
% ===========================================================================
\newpage
\section{补充分析与讨论}
\label{app:extra}

本附录汇集正文因篇幅限制未能展开的四组支撑性分析。正文各节只保留结论，
数据、逐项推理与全部对照口径均在本节完整给出。

\subsection{问题一：基准对照与灵敏度分析}

\subsubsection{基准对照：储能与光伏的价值}

为量化储能设备与光伏发电的经济价值，构造"光伏供电、缺口直接向外网购买、
电池保持不动"的基准方案：$g_t^{\mathrm{base}} = \max\{\ell_t - v_t,\,0\}$，
$c_t = d_t = 0$，$E_t \equiv 6000$，多余光伏弃用。该方案满足同样的日循环
约束，因而是原问题的一个可行解，同一口径下最优费用必然不高于它，可据此
提供独立的一致性检验。

\begin{table}[htbp]
  \centering
  \caption{不同方案的经济性对照}
  \label{tab:p1-benchmark}
  \begin{tabular}{lrrr}
    \toprule
    方案 & 全天购电量/kWh & 全天购电费/元 & 弃光电量/kWh \\
    \midrule
    无光伏无储能 & 111024.81 & 88319.36 & --- \\
    无储能       & 61789.94  & 48052.05 & 6247.96 \\
    本模型（含储能） & \textbf{59482.70} & \textbf{35126.95} & \textbf{0.00} \\
    \bottomrule
  \end{tabular}
\end{table}

由表~\ref{tab:p1-benchmark} 可见，配置储能设备后，全天购电费由 $48052.05$~元
降至 $35126.95$~元，\textbf{节省 $12925.10$~元，降幅 $26.90\%$}；较
"无光伏无储能"方案节省 $53192.41$~元（$60.23\%$）。更值得注意的是，含储能
方案的\textbf{购电量反而更少}（$59482.70$~kWh 对比 $61789.94$~kWh）：因为
储能将无储能方案中被迫弃掉的 $6247.96$~kWh 光伏盈余转移到了无光照时段使用，
从而减少了对电网的净购电需求。储能的价值不仅在于"低买高卖"的价差套利，
更在于\textbf{提升光伏本地消纳率}。优化费用 $35126.95$~元低于基准费用
$48052.05$~元，与上述可行性论证一致。

\subsubsection{灵敏度分析}

题目只给出一个"充放电效率为 $90\%$"的参数，其在充电与放电两侧的分配方式
存在解释空间：可以是单程效率（$\eta_c=\eta_d=0.9$，往返 $0.81$），也可以是
往返效率（对称折算为 $\eta_c=\eta_d=\sqrt{0.9}\approx0.9487$，往返 $0.90$）。
此外，储能的日循环水平、充放电决策的时间粒度也均可能影响结论。为此设计
六组对照口径：

\begin{table}[htbp]
  \centering
  \caption{不同建模口径下的求解结果对照}
  \label{tab:p1-sensitivity}
  \footnotesize
  \begin{tabular}{lrrr}
    \toprule
    口径 & 全天购电量/kWh & 全天购电费/元 & 弃光/kWh \\
    \midrule
    \textbf{主模型}（往返 $0.81$，$E_0=6000$，$10$ min）
      & \textbf{59482.70} & \textbf{35126.95} & 0.00 \\
    放电无损（$\eta_c=0.90$，$\eta_d=1.00$，往返 $0.90$）
      & 57580.98 & 33416.53 & 0.00 \\
    对称折算（$\eta_c=\eta_d=\sqrt{0.90}$，往返 $0.90$）
      & 57526.24 & 33801.50 & 0.00 \\
    $E_0$ 自由（仅要求 $E_0=E_{144}$，收敛至 $8550$）
      & 59482.70 & 35118.60 & 0.00 \\
    粗化为 $1$ 小时粒度
      & 59924.95 & 36771.29 & 768.44 \\
    粗化为 $4$ 小时粒度
      & 61302.10 & 43648.06 & 2293.67 \\
    \bottomrule
  \end{tabular}
  \par\vspace{2pt}
  \begin{minipage}{.92\textwidth}
    \footnotesize 注："往返"指 $\eta_c\eta_d$；除注明外均为
    $\eta_c=\eta_d=0.90$、$E_0=6000$~kWh、$10$ 分钟决策粒度。
  \end{minipage}
\end{table}

对表~\ref{tab:p1-sensitivity} 的分析可得以下结论：

\textbf{（1）效率约定影响有限，主要结论稳健。}若将 $90\%$ 的效率解释为往返
效率并对称折算（$\eta_c=\eta_d=\sqrt{0.9}$），购电费为 $33801.50$~元，较主
模型降低 $3.77\%$；若极端地假设放电无损（$\eta_c=0.90$，$\eta_d=1.00$），
购电费为 $33416.53$~元，降低 $4.87\%$。两者构成单程/往返两种解释下的费用
区间 $[33416.53,\ 35126.95]$~元，跨度不足 $5\%$。原因是放电效率越低，
储能放电时需额外多消耗的储存电量越多，越直接地削弱峰谷套利收益。\textbf{两
种口径下"谷充峰放"的策略结构与零弃光结论均不变}，说明主模型对效率参数的
分配方式不敏感。

\textbf{（2）给定的初始储电量接近最优循环水平。}若解除 $E_0 = 6000$~kWh 的
固定约束、仅要求日循环 $E_0 = E_{144}$，模型自动收敛到 $E_0 = 8550$~kWh，
购电费仅下降 $8.35$~元（$0.024\%$）。这说明附录 1 给定的 $6000$~kWh 已非常
接近成本最优的日循环工作点，无需大幅调整。该结果也验证了主模型采用
$E_0 = 6000$~kWh 的合理性：即使放开，最优解也不会显著偏离；同时它从反面
说明，把初始储电量当作可自由优化的资源只会带来不足万分之三的收益，不值得
为此偏离附录给定值。

\textbf{（3）决策粒度具有实质经济价值。}这是灵敏度分析中最显著的效应。当
充放电决策由 $10$ 分钟粗化为 $1$ 小时时，购电费上升 $4.68\%$ 并产生
$768.44$~kWh 弃光；粗化为 $4$ 小时时，购电费上升 $24.26\%$（达
$43648.06$~元），弃光量进一步增至 $2293.67$~kWh。原因是粗粒度下储能功率在
各时段内被迫恒定，无法跟踪电价与光伏的快速波动，峰谷套利的时机被显著钝化，
光伏盈余与储能充电能力也难以精确匹配，只能弃光。因此\textbf{采用 $10$ 分钟
的精细决策粒度是必要的}，这也从侧面说明题目要求以 $10$ 分钟为间隔给出购电量
是合理的。

\subsection{问题二：弃用电量的成因}

全年弃用电量达 $2463544.98$~kWh，约占同期光伏发电量
$19068890.0$~kWh 的 $12.9\%$。这一数字需要正确解读：题目未对弃电设置惩罚或
收益，因此\textbf{弃电不会带来任何额外扣款}；但它仍是日前计划与实时运行之间
偏差的\textbf{指示器}。

按逐时段口径分解全年弃电，可以看到它来自\textbf{两个性质不同}的来源。
$71.8\%$ 的弃电（$1769824.2$~kWh）发生在实际光伏大于实际负载的时段，属于
光伏富余超出储能吸纳能力——储能受 $5000$~kW 功率上限与 $10800$~kWh 容量上限
双重约束，在光伏最集中的时段无法全额吸收；风险修正又使日前计划量系统性偏高，
与高于预期的实际光伏叠加后进一步加剧了这一部分。另 $28.2\%$ 的弃电
（$693720.7$~kWh）发生在\textbf{光伏并无富余}的时段，其中 $137052.5$~kWh
甚至出现在光伏出力为零的深夜与凌晨——这部分电量不可能来自光伏，只能来自
\textbf{已按计划电价付费、但因实际负载低于预期而用不掉的计划购电量}。按各
时段电价核算，这部分被弃掉的已购电量当初共付出 $466267$~元，约占全年计划
购电费的 $3.54\%$。换言之，弃电并非严格"零成本"：其中约三成是"按 $5$ 倍
电价买保险"所预付的保费中未能出险的部分。

从经济学角度看，这一取舍是\textbf{经过权衡的最优选择}：减少弃电最直接的
办法是降低计划购电量，但那会同时削弱对紧急购电的防护。本文在固定
$\rho=1$、$\lambda=0$ 下重跑了 $\alpha\in\{0.50,0.70,0.80,0.90\}$ 四点的
敏感性验证（逐点见支撑材料 \texttt{p2\_alpha\_sensitivity.json}）：$\alpha$
由 $0.80$ 降至 $0.50$ 时，弃电量从 $2628.6$~MWh 降至 $1493.7$~MWh，但紧急
购电量由 $177.9$~MWh 猛增至 $636.1$~MWh，\textbf{总费用反而从 $1403.8$ 万元
升至 $1508.8$ 万元}；反向升至 $0.90$ 时，紧急购电量降至 $86.2$~MWh，弃电量
却升至 $3369.1$~MWh，总费用回升至 $1421.6$ 万元；$\alpha=0.70$ 为
$1410.8$ 万元。四点中 $\alpha=0.80$ 取到最小值且两侧点均高于它，
"以适量弃电换取紧急购电防护"的取舍是可验证的。（这四点均固定 $\rho=1$；
主策略因滚动重标定在 $8$ 个窗口选到不同参数，全年 $1393.97$ 万元还略低于该
网格最优点。）

图~\ref{fig:p2-daily} 给出紧急购电量的季节分布与逐日费用构成，其中还暴露了一个
对风险管理很有意义的事实：\textbf{紧急购电的负担在时间上高度集中}。
$334$ 天中有 $143$ 天发生紧急购电，但其中位日仅 $317.1$~kWh；而最严重的
$6$ 月 $1$ 日单日紧急购电 $10545.11$~kWh、费用 $56753.15$~元，\textbf{一天就
占去全年紧急购电费的 $7.5\%$}，其后紧接的 $6$ 月 $2$ 日（$7244.83$~kWh、
$42449.26$~元）与 $12$ 月 $2$ 日、$7$ 月 $6$ 日同样显著高出常态。
这说明全年紧急购电费并非由 $143$ 个"小缺口"平摊，而是由少数几个
\textbf{连续多日的高压时段}主导——$6$ 月初的连日高负荷恰逢光伏出力波动，
使储能在数日内持续被压至下限，缺口只能在 $5$ 倍电价下补足。

这一分布形态具有直接的运营含义：若把有限的额外储能容量或预测投入用于降低
\textbf{极端日}的缺口，其边际收益将远高于用于改善 $150$~kWh 级别的日常误差。

\subsection{问题二：模型的边界}

需要明确本文模型的适用边界。本文的"风险余量 + 日前 MILP + 因果执行规则"
是一条便于实现和解释的\textbf{近似策略}，其优劣完全由按时间顺序的实际运行
回放判定，\textbf{并不等同于已求出所有随机控制策略中的全局最优解}。当
$\lambda=0$ 时，日前模型的目标中不含任何对未来的价值评估，因此它\textbf{}
不会主动为次日保留电量，其合理性由跨日连续状态的自然约束保证，而非由模型
内的随机优化保证。更严格的改进方向是把执行策略显式建模为非预见性约束下的
场景优化（所有场景共用同一份 $0{:}00$ 购电计划，且观测历史相同的场景必须
作出相同动作），或采用因果滚动优化，每步只用当前观测值重排后续储能轨迹。
就本题给定的数据规模与 $5000$~kW 功率上限而言，现行方案已经取得了较"无储能"
基准 $24.73\%$ 的费用节约，是否值得进一步引入上述模型，应依据其在历史回放
中的实际改善与计算代价决定。

\subsection{问题三：模型的边界}

\textbf{第一，交付块策略是一种结算口径的选择，不是题目的唯一解。}本文限定
每个六小时交付块只在块首调整一次，从而让每个时段只与自己的 $g_t^0$ 结算一次。
若允许在 $6{:}00$、$12{:}00$ 反复修订同一个傍晚时段，就必须另行声明是"最终
净偏差结算一次"还是"逐次变更分别收费"，两者会得到不同的费用。本文的结论
首先评价的是这一基础策略。

\textbf{第二，本文求的是可实现的滚动近似策略，而不是未知未来下的全局最优
解。}每个决策时点只使用当时已经发布的信息，没有对未来尚未发布的预报作多阶段
随机优化，因此 $14{,}977{,}809$~元（对照）与 $13{,}418{,}626$~元（主策略）都是
\textbf{相对于信息集可达到的费用}，而不是全知情形下的下界。

\textbf{第三，参数标定是 $31$ 天窗口内的近视最优。}网格搜索的目标是 2025 年
$1$ 月的实际总费用，它并不保证全年最优。正文的标定小节已经用样本外检验量化了
这一风险：1 月标定面上的最小点
$(\alpha,\rho,W)=(0.6,0.35,14)$ 在 1 月上比采用的 $(0.6,0.5,28)$ 省
$9{,}725$~元，但在评价区间上\textbf{反而多花 $43{,}213$~元}。这是"不许用未来
数据选参"这一硬约束的必然代价。因此表中各方式的费用差异应当理解为\textbf{同一
规则、同一组冻结参数下的可比结果}，而不是各方式的最优可能值。

\textbf{第四，计费口径的歧义会影响调减的方向与账本。}本文主口径把调减理解为
“退回原价、只收 $50\%$ 违约费”，此时调减每度净得 $0.5p$，最优解中会
真实出现调减（全年 $775{,}660.56$~kWh）。若按“调减仍付原价”理解，调减
永不优于不调，最优解不会真正调减。本文把后一种读法作为独立的敏感性检查一并报告
（表~\ref{tab:p3-billing}），两臂在全年总费用上相差 $1.49\%$；\textbf{本文不对
哪一种读法作对错判断}，只报告差异的量级。

\textbf{第五，时段粒度与记账时点。}所有结果建立在 $10$ 分钟离散记账、且执行层
可对每段实际负载与光伏作\textbf{段内实时测量}（理想即时响应、无控制滞后）的
假设上；实际运行中若观测与调度存在滞后，需要在执行层加入相应的安全裕量。

\subsection{问题四：模型的边界}
\label{app:p4-limit}

本文的模型有若干明确的边界，说明如下，以免读者把结论推广到未被验证的范围。

\textbf{价格预测的能力边界。}附件 4 与附件 1 的\textbf{年平均完全相同}、
逐时段均值最大偏差仅 $5.15\times10^{-5}$~元/kWh，二者之差是一个与附件 1 曲线
无关的残差；该残差中\textbf{星期结构解释 $41.7\%$ 的方差}，其余部分本文按
不可预测处理。这意味着：\textbf{如果现实中价格包含真正的随机成分（如机组
非计划停运、极端天气导致的负荷尖峰），本文的预测器无法预知，只能靠风险余量
对冲}。$\Delta J$ 与图~\ref{fig:p4-strategy} 中"未来价格已知参考"之间的差距，
正是这一不可预测部分的价格。

\textbf{风险处理是"分位数余量"而非"场景优化"。}本文没有构造价格—负载的联合
场景集，而是把风险净负荷 $\widetilde N_t=\widehat N_t+Q_\alpha(\varepsilon_t)$
作为计划依据；这一做法要求误差分布在一段时间内平稳，且只保证\textbf{边际}
上的安全，不保证在极端联合情形下仍然最优。第 \ref{sec:p4-risk} 节已经说明，
价格与误差的协方差为正，忽略它会\textbf{低估}紧急购电风险，本文通过价格加权
分位数把它显式计入，但没有进一步做鲁棒优化。

\textbf{计费口径的歧义。}题目没有说明"计划量高于调整量的部分按 $50\%$ 退"
是\textbf{退回原价}还是\textbf{只按半价结算}，也没有说明调减是否涉及退款。
主口径取"不退款、只追加"，并对退款口径单独做了敏感性回放（见附录
第 \ref{app:p4-refund} 节）；两种口径下 4-3 与 4-2 的相对关系不变，说明主结论
不依赖这一选择。

\textbf{参数搜索网格截断了 4-3 的最优解。}$4$-$3$ 的 $(\alpha,\rho)$ 在 $8$ 个
标定窗口里全部停在网格左下角，且网格内费用沿两个方向单调递增，有 $4$ 个窗口
胜出者与次优者仅差 $0.19$--$25.12$ 元，说明\textbf{真实最优多半落在网格之外}
（$\rho<0.5$ 甚至 $\rho=0$ 从未被试过），落在角上时求解器返回哪个极点带有
任意性。本文没有扩网格重跑（4-3 每轮标定需解 $1336$ 次 MILP），故 \textbf{4-3
的 $\alpha,\rho$ 取值不具可解释性}；相对地，4-2 的网格内费用\textbf{不单调}，
$8$ 个窗口选中不同的 $(\alpha,\rho)$ 且次优差达 $26.9$--$3112.5$ 元。

\textbf{未建模的部分。}本文不含储能寿命与循环次数的成本、不含需求侧响应、
不含电网侧的安全约束（如联络线容量、功率因数），也不考虑电价预测误差对
\textbf{次日以后}的影响。这些都属于模型之外的因素，把它们纳入会改变具体
数值，但不改变"按当时信息决策、按实际价格结算"这一框架本身的结论。
